\sectiontitle{Q2. Functional-Block Literature and Development Map}

{\fontsize{9.6}{11.35}\selectfont
The assignment requires more than a general demonstration that the concept can
work: it asks which functional blocks already have reasonable solutions that
can be leveraged, which blocks still require project-specific integration, and
where the strongest technical uncertainty remains. The comparison below treats
published prosthetic-interface systems and printed-sensor studies as
\textbf{similar or competing research solutions}. It does not claim that each
paper reproduces the complete liner described in this report.\par}

\hypertarget{q2-functional-blocks}{}
\qtwosubsection{2.1 Similar and competitive solutions by functional block}

\begingroup
\fontsize{8.2}{9.65}\selectfont
\renewcommand{\arraystretch}{1.09}
\begin{tabularx}{\linewidth}{>{\raggedright\arraybackslash\bfseries\color{AUBRed}}p{2.45cm}Y Y Y}
\rowcolor{AUBRed}
\color{white}\textbf{Functional block} &
\color{white}\textbf{Closest published solution} &
\color{white}\textbf{Reasonable solution to leverage} &
\color{white}\textbf{Project-specific integration work} \\
Pressure and directional shear & Four partially overlapping printed
capacitors separate normal and signed shear response
\cite{albrecht2018conf,albrecht2019}; textile arrays demonstrate distributed
socket-pressure measurement \cite{tabor2021}. & Reuse the four-sector cell
principle, differential capacitance readout, and calibrated directional loading
method. & Print and protect the cell on the 5.2 mm PDMS slab, quantify
cross-axis error, and preserve its physical coordinate in the multimodal map. \\
\rowcolor{AUBPale}
Temperature & A printed silver meander provides a repeatable
resistance--temperature relationship on a flexible substrate \cite{liew2020}. &
Reuse the serpentine resistive geometry and temperature-calibration framework. &
Separate temperature response from bending, compression, conductor strain, and
self-heating after direct printing on PDMS. \\
Humidity and sweat route & An all-inkjet-printed interdigitated capacitive
sensor measures relative humidity through a hygroscopic layer
\cite{aitmammar2020}. & Reuse the IDE geometry, capacitance measurement, and
vapour-step characterization. & Provide vapour access while limiting liquid
sweat, leakage, saturation, slow recovery, and temperature dependence inside
the liner. \\
\rowcolor{AUBPale}
Printed PDMS conductors & Plasma-assisted AgNP printing and partially embedded
silver traces establish two direct-on-PDMS process routes
\cite{albrecht2018stretch,sun2016}. & Use surface treatment, wide-trace process
windows, and bending/strain characterization as starting methods. & Determine
the actual print window, swelling, adhesion, geometry, curing, and durability
for the project's ink, printer, sensor artwork, and molded PDMS. \\
Sensorized liner integration & Embedded temperature/humidity sensing and a
wireless multimodal skin-interface platform show that prosthetic-liner sensing
can be distributed and synchronized \cite{paterno2020,q3kwak2020}. & Reuse
staged integration, channel identification, synchronized acquisition, and
post-assembly calibration. & Combine all four printed modalities in thin slabs,
route connections through a molded liner, and retain calibration under
curvature, pre-strain, donning, and repeated loading. \\
\rowcolor{AUBPale}
Spatial interpretation & Socket pressure maps and movement/pistoning systems
demonstrate spatial display and independent motion ground truth
\cite{q3turner2021,q3yu2024,noll2018}. & Reuse coordinate registration,
known-location stimulation, interpolation checks, and movement reference
measurements. & Fuse pressure, shear, temperature, and humidity without hiding
modality-specific error; report localization error before interpreting a map as
discomfort risk. \\
Skin interface and bench model & Hydrogel--elastomer bonding, hydrogel sweat
collection, and prosthetic socket phantoms provide separate material and test
precedents \cite{yuk2016,lu2025,quinlan2020,phillips2025}. & Reuse established
bonding concepts, controlled moisture comparisons, and mechanically repeatable
phantom loading. & Compare bare and patterned skin-facing structures without
blocking or redistributing the signals, then transfer passing slabs to a curved
liner geometry. \\
\end{tabularx}
\endgroup
\sourceanchor{Functional-block comparison synthesized from printed-sensor,
sensorized-liner, visualization, interface-material, and prosthetic test-rig
literature. Each project-specific statement is bounded to work that remains to
be demonstrated on this design.}

\vspace{0.5em}
\evidencebox{Novelty boundary:}{The individual sensor principles, PDMS printing
methods, hydrogel bonding concepts, and spatial displays are \textbf{not treated
as independently novel}. The candidate project contribution is their
coordinate-preserving integration as directly printed pressure, shear,
temperature, and humidity blocks on modular 5.2 mm PDMS slabs, followed by
multimodal calibration and liner-level mapping. A formal novelty claim would
still require a dedicated patent and systematic prior-art search.}

\vspace{0.45em}
\begin{center}
\begin{tikzpicture}[>=stealth,node distance=2.6mm]
  \tikzstyle{mapstep}=[draw=AUBLine,fill=AUBPale,rounded corners=1pt,
    minimum height=7.5mm,text width=34mm,align=center,
    font=\fontsize{6.9}{7.8}\selectfont]
  \node[mapstep] (m1) {established sensing\\and printing blocks};
  \node[mapstep,right=of m1] (m2) {project PDMS slab\\and CAD coordinates};
  \node[mapstep,right=of m2] (m3) {multimodal calibration\\and protection};
  \node[mapstep,right=of m3] (m4) {curved liner and\\spatial interpretation};
  \draw[->,thick,color=AUBRed] (m1)--(m2);
  \draw[->,thick,color=AUBRed] (m2)--(m3);
  \draw[->,thick,color=AUBRed] (m3)--(m4);
\end{tikzpicture}
\end{center}
{\fontsize{7.5}{8.7}\selectfont\color{AUBGray}
\textbf{Figure Q2-F1. Functional-block strategy.} Existing solutions are used
as engineering starting points; the work concentrates on the interfaces between
blocks and on the evidence required after each mechanical boundary changes.\par}

\clearpage
\sectiontitle{Q2. Work Completed and Current Build}

\hypertarget{q2-current-build}{}
\qtwosubsection{2.2 What has been made and what is active now}

{\fontsize{9.65}{11.4}\selectfont
Work already completed includes printing the full sensor artwork on a large PET
sheet, molding the 5.2 mm PDMS slab, and developing the Fusion/AutoCAD sensor
layouts and slab renders. PET was the available first-stage printing substrate,
not a candidate liner material. The active build transfers the same functional
blocks to the molded PDMS specimen, whose CAD envelope is approximately
\textbf{56 mm $\times$ 55 mm $\times$ 5.2 mm}.\par}

\begingroup
\fontsize{8.85}{10.35}\selectfont
\renewcommand{\arraystretch}{1.14}
\begin{tabularx}{\linewidth}{>{\raggedright\arraybackslash\bfseries\color{AUBRed}}p{3.0cm}Y Y}
\rowcolor{AUBRed}
\color{white}\textbf{Artifact or workstream} &
\color{white}\textbf{Present state} &
\color{white}\textbf{Role in the development sequence} \\
Large PET process print & The complete sensor artwork has been printed on the
available large PET sheet. & Records artwork transfer, printer operation,
trace routing, connection access, and first-stage electrical inspection. It
does not determine the liner substrate. \\
\rowcolor{AUBPale}
Molded PDMS slab & A nominal 5.2 mm PDMS specimen has been formed for direct
printing. & Provides the intended flexible substrate for recording cure state,
surface preparation, printed geometry, swelling, adhesion, and sensor response. \\
Sensor artwork & Pressure/shear capacitor geometry, resistive temperature
meander, capacitive humidity IDE, conductors, and pads are available. & Defines
the four established sensing blocks transferred to the PDMS print and later
checked channel by channel. \\
\rowcolor{AUBPale}
CAD slab and placement design & Fusion/AutoCAD renders define the slab envelope,
relative sensor arrangement, central enclosure region, and conductor exits. &
Supplies the coordinate reference for the printed build; a final dimensioned
export will distinguish nominal from measured dimensions. \\
Acquisition and mapping algorithm & Threshold logic and spatial-map generation
are under development. & Receives calibrated channels only after sensor identity,
units, time synchronization, and CAD coordinates are documented; algorithmic
validation is addressed later in the report. \\
\end{tabularx}
\endgroup
\sourceanchor{Project status derived from the completed PET print, molded PDMS
specimen, supplied Fusion/AutoCAD geometry, slab renders, and stated software
work. The literature supplies transfer methods and measurement precedents, not
evidence that these project artifacts have already been calibrated
\cite{albrecht2019,liew2020,aitmammar2020,albrecht2018stretch}.}

\evidencebox{Drawing requirement:}{The next CAD export must label the coordinate
origin, active footprints, conductor widths, connection pads, edge clearances,
enclosure sections, and verified versus nominal dimensions. Perspective renders
must not be used as dimensional evidence. The supplied render and placement
blueprint are consolidated with the dimensional requirements and published
liner-thickness context in
\hyperlink{q4-implementation-geometry}{\textcolor{AUBRed}{\textbf{Q4 Section 4.2}}}.}

\vspace{0.55em}
\begin{minipage}[t]{0.29\textwidth}
  \vspace{0pt}\centering
  \includegraphics[width=\linewidth,height=66mm,keepaspectratio,
    trim=285 365 285 18,clip]
    {assets/project/pet_sensor_process_print.png}
\end{minipage}\hfill
\begin{minipage}[t]{0.685\textwidth}
  \vspace{0pt}
  {\fontsize{8.55}{10.1}\selectfont\color{AUBGray}\raggedright
  \textbf{Figure Q2-W1. Sensors printed on PET during the first-stage
  fabrication trial.} The visible layouts correspond, from top to bottom, to
  the multi-sector pressure/directional-shear electrode pattern, the resistive
  temperature meander, and the dense capacitive humidity IDE. This project
  photograph documents successful artwork transfer and printing on the
  available PET process substrate.\par
  \vspace{0.45em}
  \textbf{Evidence boundary.} The photograph does not establish ink
  conductivity, calibration, sensitivity, channel isolation, or performance
  after direct printing on molded PDMS; those require electrical inspection
  and the modality-specific validation described later in the report.\par}
\end{minipage}

\clearpage
\sectiontitle{Q2. Functional-Block Literature---Pressure and Shear}

{\fontsize{9.7}{11.5}\selectfont
Across the reviewed force-sensing work, the consensus is that local normal and
directional shear information can be recovered from differential sensor
responses. The methods differ in substrate, enclosure, electrode geometry,
readout, and loading range, so their numerical performance is \emph{source
performance}. The established sensing principle can be reused, but the values
cannot be transferred to the project's PDMS stack without recalibration.\par}

\hypertarget{q2-pressure-shear}{}
\qtwosubsection{2.3 Pressure and directional shear}

\begin{minipage}[t]{0.38\textwidth}
  \vspace{0pt}
  {\fontsize{9.2}{10.8}\selectfont
  Albrecht \textit{et al.} used four partially overlapping printed capacitors
  to distinguish normal loading from directional shear
  \cite{albrecht2018conf,albrecht2019}. The four capacitances, C1--C4, respond
  differently with force direction. The journal device reported an
  approximately linear response over 0.1--8 N, with 5.2 fF/N normal-force and
  13.1 fF/N shear-force sensitivity. Those values belong to the published
  device and cannot be assigned to the present slab before calibration.\par

  \vspace{0.55em}
  \textbf{Transfer requirement.} Apply known normal loads and signed shear in
  both axes. Report sensitivity, linearity, hysteresis, cross-axis response,
  repeatability, and spatial localization for every printed cell.}
\end{minipage}\hfill
\begin{minipage}[t]{0.59\textwidth}
  \vspace{0pt}
  \centering
  \includegraphics[width=\linewidth,height=72mm,keepaspectratio]
    {assets/q2/sensor_literature/pressure_shear_sensor_evidence.png}\par
  \vspace{0.2em}
  {\fontsize{8.15}{9.4}\selectfont\color{AUBGray}\raggedright
  \textbf{Figure Q2-S1. Pressure/shear evidence plate.} Project artwork and the
  published four-capacitor force-cell principle reported by Albrecht
  \textit{et al.} \cite{albrecht2019}.\par}
\end{minipage}

\vspace{0.75em}
\begingroup
\fontsize{8.9}{10.45}\selectfont
\renewcommand{\arraystretch}{1.14}
\begin{tabularx}{\linewidth}{>{\raggedright\arraybackslash\bfseries\color{AUBRed}}p{2.8cm}Y Y}
\rowcolor{AUBRed}
\color{white}\textbf{Comparison dimension} &
\color{white}\textbf{What existing solutions establish} &
\color{white}\textbf{Adaptation in the current project} \\
Normal versus shear & Four capacitor sectors change differently with force
direction. & A calibrated force matrix that reconstructs normal load and signed
shear without unacceptable cross-talk. \\
\rowcolor{AUBPale}
Sensor placement & Local force cells can resolve the load applied at their
footprint. & Coverage and missed-hotspot tests using known locations between
and around the current sensor sites. \\
Durability & Published flexible devices report function under their own
material stacks. & Repeated compression, bidirectional shear, bending, and
recalibration on the molded slab. \\
\end{tabularx}
\endgroup
\sourceanchor{Literature-derived mechanism summary from the four-capacitor
printed force-cell studies \cite{albrecht2018conf,albrecht2019}; the right-hand
column defines the adaptation from the published solution to this project.}

\vspace{0.5em}
\evidencebox{Reuse decision:}{Adopt the four-sector capacitive architecture and
directional calibration logic. Develop only the substrate-specific printing,
PDMS enclosure, strain/cross-axis compensation, and coordinate registration
needed to make that established block operate inside the multimodal slab.}

\clearpage
\sectiontitle{Q2. Functional-Block Literature---Temperature and Humidity}

\hypertarget{q2-temperature}{}
\qtwosubsection{2.4 Printed temperature sensing}

\begin{minipage}[t]{0.38\textwidth}
  \vspace{0pt}
  {\fontsize{9.15}{10.75}\selectfont
  Liew \textit{et al.} printed serpentine silver resistors on PET and observed
  a linear resistance--temperature relationship from 30 to 100~$^\circ$C
  \cite{liew2020}. Their reported design was approximately
  6.9 mm $\times$ 13.2 mm, with 0.5 mm tracks and 0.4 mm gaps. Resistance also
  depends on printed geometry and mechanical strain, so the present meander
  requires temperature calibration while the slab is unloaded and under
  representative bending and compression.\par

  \vspace{0.45em}
  \textbf{Required outputs.} Temperature coefficient, offset, response time,
  repeatability, strain error, self-heating check, and drift after cycling.}
\end{minipage}\hfill
\begin{minipage}[t]{0.59\textwidth}
  \vspace{0pt}
  \centering
  \includegraphics[width=\linewidth,height=54mm,keepaspectratio]
    {assets/q2/sensor_literature/temperature_sensor_evidence.png}\par
  \vspace{0.15em}
  {\fontsize{8.1}{9.35}\selectfont\color{AUBGray}\raggedright
  \textbf{Figure Q2-S2. Temperature evidence plate.} Project meander artwork
  and the published resistive temperature-sensor principle from Liew
  \textit{et al.} \cite{liew2020}.\par}
\end{minipage}

\vspace{0.85em}
\hypertarget{q2-humidity}{}
\qtwosubsection{2.4.1 Printed humidity sensing}

\begin{minipage}[t]{0.38\textwidth}
  \vspace{0pt}
  {\fontsize{9.15}{10.75}\selectfont
  A\"{\i}t-Mammar \textit{et al.} demonstrated an all-inkjet-printed
  capacitive relative-humidity sensor using interdigitated silver electrodes
  and a hygroscopic sensing layer \cite{aitmammar2020}. The supplied plate
  connects the project IDE to the published geometry and micrograph. The paper
  supports the capacitive mechanism, but it does not prove liquid-sweat
  resistance or the response of this device inside a prosthetic liner.\par

  \vspace{0.45em}
  \textbf{Required outputs.} Vapour-response curve, temperature compensation,
  response and recovery time, saturation, liquid-exposure survival, hysteresis,
  drift, and repeatability after drying.}
\end{minipage}\hfill
\begin{minipage}[t]{0.59\textwidth}
  \vspace{0pt}
  \centering
  \includegraphics[width=\linewidth,height=54mm,keepaspectratio]
    {assets/q2/sensor_literature/humidity_sensor_evidence.png}\par
  \vspace{0.15em}
  {\fontsize{8.1}{9.35}\selectfont\color{AUBGray}\raggedright
  \textbf{Figure Q2-S3. Humidity evidence plate.} Project IDE artwork and
  published capacitive humidity-sensor evidence from A\"{\i}t-Mammar
  \textit{et al.} \cite{aitmammar2020}.\par}
\end{minipage}

\vspace{0.7em}
\evidencebox{Consensus, difference, and gap:}{Printed resistive meanders and
capacitive IDEs are established solutions for temperature and humidity,
respectively. Their functional layers and measurands differ: the temperature
trace must reject strain effects, whereas the humidity IDE requires a controlled
vapour route and liquid-sweat protection. The project therefore reuses the two
sensor principles while developing their PDMS-specific calibration, protection,
cross-sensitivity control, and common coordinate mapping.}

\clearpage
\sectiontitle{Q2. Direct-on-PDMS Printing Literature}

\hypertarget{q2-pdms-printing}{}
\qtwosubsection{2.5 Process routes, transfer limits, and project controls}

{\fontsize{9.55}{11.25}\selectfont
Two direct-printing studies provide process solutions that can be leveraged
when transferring the artwork from the successful PET print. Albrecht
\textit{et al.} used oxygen-plasma-treated silicone and inkjet-printed AgNP
conductors; Sun \textit{et al.} printed a silver precursor onto precured PDMS
to create a partially embedded, semi-wrapped conductor
\cite{albrecht2018stretch,sun2016}. The routes agree that PDMS surface and
mechanical behaviour must be designed into the print process, but they use
different chemistries and do not supply calibration for the present sensor
stack.\par}

\begin{minipage}[t]{0.485\textwidth}
  \vspace{0pt}
  \centering
  \includegraphics[width=\linewidth,height=49mm,keepaspectratio]
    {paper_sources/figures_selected/albrecht2018_fig2_pdms_pet_height_profile.png}\par
  \vspace{0.15em}
  {\fontsize{7.95}{9.15}\selectfont\color{AUBGray}\raggedright
  \textbf{Figure Q2-P1. Same-printer height profiles.} Approximate average
  heights were 3.8~$\mu$m on PDMS and 0.6~$\mu$m on PET. The apparent PDMS
  height includes substrate swelling, not only deposited silver
  \cite{albrecht2018stretch}. CC BY 4.0.\par}
\end{minipage}\hfill
\begin{minipage}[t]{0.485\textwidth}
  \vspace{0pt}
  \centering
  \includegraphics[width=\linewidth,height=49mm,keepaspectratio]
    {paper_sources/figures_selected/sun2016_fig6_bending_resistance.png}\par
  \vspace{0.15em}
  {\fontsize{7.95}{9.15}\selectfont\color{AUBGray}\raggedright
  \textbf{Figure Q2-P2. Bending comparison.} Sun \textit{et al.} reported a
  semi-wrapped PDMS conductor with substantially better resistance retention
  during bending than the paper's PET comparison, including repeated bending
  shown by the normalized resistance \cite{sun2016}. CC BY 4.0.\par}
\end{minipage}

\vspace{0.65em}
\begingroup
\fontsize{8.55}{10.05}\selectfont
\renewcommand{\arraystretch}{1.11}
\begin{tabularx}{\linewidth}{>{\raggedright\arraybackslash\bfseries\color{AUBRed}}p{2.5cm}Y Y}
\rowcolor{AUBRed}
\color{white}\textbf{Literature finding} &
\color{white}\textbf{What it means} &
\color{white}\textbf{Project control} \\
Surface wetting & Untreated PDMS is difficult to wet. In the Albrecht process,
oxygen plasma enabled continuous lines and printing occurred within 30 min of
treatment. & Record treatment power/time, contact angle if available, and the
treatment-to-print interval; do not assume the paper's recipe transfers
unchanged. \\
\rowcolor{AUBPale}
Print fidelity & The paper reports satellite drops from printer inaccuracy and
local failures in narrow lines; 1--2 mm lines were selected for reliable
stretchable conductors in that reported process. & Measure printed width,
edge roughness, gaps, shorts, open traces, and yield for the project's actual
sensor geometries. \\
Solvent interaction & The same deposited ink produced a much larger apparent
profile on PDMS. The authors attributed most of the difference to solvent
absorption and swelling within roughly the upper 15~$\mu$m. & Profile untreated,
wet-printed, dried, and equilibrated specimens; separate conductor thickness
from substrate deformation. \\
\rowcolor{AUBPale}
Mechanical survival & Printed silver remained conductive after more than
10,000 reported 20\% stretch cycles, while the semi-wrapped approach maintained
resistance through repeated bending. & Test the present ink and geometry under
compression, bending, and shear; record in-cycle resistance, recovery time,
drift, cracking, and adhesion. \\
Measurement accuracy & These papers demonstrate conductors and process
strategies, not accurate pressure, shear, temperature, and humidity sensing in
this slab. & Establish modality-specific calibration, error, hysteresis,
repeatability, cross-sensitivity, drift, and localization after printing. \\
\end{tabularx}
\endgroup
\sourceanchor{Literature-derived PDMS printing constraints and project controls
from Albrecht \textit{et al.} and Sun \textit{et al.}
\cite{albrecht2018stretch,sun2016}.}

\vspace{0.5em}
\evidencebox{Method selected from the literature:}{Use direct printing on the
molded PDMS slab, beginning with documented surface treatment and a print-window
study. PET remains process evidence only. The Albrecht and Sun routes supply
starting controls for wetting, trace width, swelling, adhesion, bending, and
cycling; the final recipe and all sensor calibration are established on the
project's own ink--printer--PDMS combination.}

\clearpage
\sectiontitle{Q2. Phase 3: Hydrogel Skin-Interface Study}

\hypertarget{q2-hydrogel-study}{}
\qtwosubsection{2.6 Literature evidence and the sensor-interference problem}

{\fontsize{9.55}{11.25}\selectfont
This is an \textbf{evidence-led development stage following successful printing
and calibration on the molded PDMS slab}. A continuous
hydrogel between skin and sensors can improve softness or hold moisture, but it
can also alter every quantity being measured. It is therefore treated
as a separate experimental variable rather than assumed to be a neutral comfort
layer.\par}

\begin{minipage}[t]{0.485\textwidth}
  \vspace{0pt}
  \centering
  \includegraphics[width=\linewidth,height=54mm,keepaspectratio]
    {paper_sources/figures_selected/yuk_fig1_hydrogel_elastomer_fabrication.png}\par
  \vspace{0.15em}
  {\fontsize{7.95}{9.15}\selectfont\color{AUBGray}\raggedright
  \textbf{Figure Q2-H1. Hydrogel--elastomer bonding precedent.} Yuk
  \textit{et al.} preserved microstructures while covalently bonding tough
  hydrogels to elastomers including Sylgard 184 PDMS \cite{yuk2016}. The figure
  supports a bonding method, not prosthetic-liner sensor accuracy. CC BY 4.0.\par}
\end{minipage}\hfill
\begin{minipage}[t]{0.485\textwidth}
  \vspace{0pt}
  \centering
  \includegraphics[width=\linewidth,height=54mm,keepaspectratio]
    {paper_sources/figures_selected/lu_fig2_hydrogel_thickness_sweat_absorption.png}\par
  \vspace{0.15em}
  {\fontsize{7.95}{9.15}\selectfont\color{AUBGray}\raggedright
  \textbf{Figure Q2-H2. Hydrogel thickness and sweat uptake.} Lu
  \textit{et al.} compared agarose--glycerol sheets of different thicknesses;
  water uptake and equilibration changed with thickness \cite{lu2025}. This was
  a wearable patch, not a load-bearing prosthetic liner. CC BY 4.0.\par}
\end{minipage}

\vspace{0.65em}
\begingroup
\fontsize{8.65}{10.15}\selectfont
\renewcommand{\arraystretch}{1.12}
\begin{tabularx}{\linewidth}{>{\raggedright\arraybackslash\bfseries\color{AUBRed}}p{2.65cm}Y Y}
\rowcolor{AUBRed}
\color{white}\textbf{Evidence theme} &
\color{white}\textbf{What the literature supports} &
\color{white}\textbf{What remains unproven here} \\
Bonding to PDMS & Robust covalent interfaces and preserved microstructures are
possible with specific surface chemistry and UV processing
\cite{yuk2016}. & Compatibility with the printed silver, sensor films, curing
history, cleaning, skin contact, and the project's fabrication equipment. \\
\rowcolor{AUBPale}
Moisture uptake & Hydrogel sheets can collect sweat, and uptake depends on
thickness and time \cite{lu2025}. & Whether uptake prevents pooling or instead
stores sweat beside the skin; release, drying, hygiene, and repeat-cycle
behaviour are unknown. \\
Sensor placement & A personalised liner study embedded temperature and humidity
sensors in direct skin contact and found that integration geometry is central
to interpretation \cite{paterno2020}. & The response after inserting a
water-rich intermediate layer cannot be inferred from direct-contact data. \\
\rowcolor{AUBPale}
Material mechanics & Swollen hydrogels are compliant and can change dimensions;
Yuk \textit{et al.} treated interface toughness and microstructure preservation
as design requirements \cite{yuk2016}. & Pressure transmission, shear transfer,
friction, swelling-induced preload, spatial blur, and cyclic durability in this
slab. \\
\end{tabularx}
\endgroup
\sourceanchor{Literature synthesis of hydrogel--elastomer bonding, sweat uptake,
and direct-contact liner sensing \cite{yuk2016,lu2025,paterno2020}.}

\vspace{0.5em}
\evidencebox{Literature conclusion:}{The literature supports a carefully gated
hydrogel experiment, not immediate adoption of a continuous skin-facing layer.
No reviewed paper validates this hydrogel arrangement over the complete
combination of
pressure, directional shear, temperature, and humidity sensors in a prosthetic
liner.}

\clearpage
\sectiontitle{Q2. Phase 3: Hydrogel Study---Continued}

\hypertarget{q2-hydrogel-interactions}{}
\qtwosubsection{2.6.1 Sensor interactions introduced by the hydrogel}

\begingroup
\fontsize{8.8}{10.35}\selectfont
\renewcommand{\arraystretch}{1.13}
\begin{tabularx}{\linewidth}{>{\raggedright\arraybackslash\bfseries\color{AUBRed}}p{2.1cm}Y Y}
\rowcolor{AUBRed}
\color{white}\textbf{Modality} &
\color{white}\textbf{Likely interference mechanism} &
\color{white}\textbf{Required paired test} \\
Humidity & Water absorption and diffusion can buffer a rapid humidity change,
delay recovery, create hysteresis, or saturate the environment near the IDE.
Liquid or ionic pathways may also increase leakage. & Compare H0 bare PDMS, H1
islands, the H2 windowed sheet, and H3 continuous hydrogel under vapour steps,
controlled sweat volume, drying, and repeated wet/dry cycles. \\
\rowcolor{AUBPale}
Temperature & Added thermal mass, thermal resistance, absorbed water, and
evaporation can change both the steady temperature and response time at the
skin interface. & Apply identical temperature steps and heat fluxes; report
offset, time constant, spatial smoothing, and recovery for dry and wet states. \\
Pressure & A soft swollen layer can redistribute normal stress, change contact
area, attenuate peaks, introduce preload, and add viscoelastic hysteresis. & Use
known loads and footprints; compare sensitivity, linearity, hysteresis, drift,
peak attenuation, and hotspot localization before and after swelling. \\
\rowcolor{AUBPale}
Directional shear & Friction, adhesion, slip, and strain transfer change with
water content. A continuous layer may reduce, rotate, or spread the shear
reaching the four-capacitor cell. & Apply signed shear in both axes at controlled
normal preload; quantify vector error, cross-axis coupling, slip, and spatial
blur in dry and wet conditions. \\
Printed silver & Water, ions, swelling strain, and bonding chemistry may affect
insulation, adhesion, continuity, corrosion, or contact resistance. & Measure
leakage, resistance, insulation, adhesion, cracking, and recovery through cyclic
loading and wet/dry exposure. \\
\end{tabularx}
\endgroup
\sourceanchor{Sensor-interference analysis derived from the material
transport and bonding evidence reviewed in Section 2.6
\cite{yuk2016,lu2025,paterno2020}.}

\vspace{0.25em}
\begin{center}
\begin{tikzpicture}[x=1cm,y=1cm]
  \foreach \x/\lab/\col in {0/H0: bare PDMS/AUBGray,
                             4.0/H1: hydrogel islands/AUBGreen,
                             8.0/H2: windowed sheet/AUBYellow,
                             12.0/H3: continuous layer/AUBCoral}{
    \fill[AUBPale] (\x,0) rectangle ++(3.2,0.38);
    \draw[AUBGray] (\x,0) rectangle ++(3.2,0.38);
    \draw[AUBRed,very thick] (\x+0.55,0.18)--(\x+2.65,0.18);
    \node[below,font=\fontsize{6.8}{7.7}\selectfont,align=center]
      at (\x+1.6,-0.08) {\lab};
  }
  \fill[AUBGreen!60] (4.25,0.42) rectangle (4.9,0.65);
  \fill[AUBGreen!60] (6.3,0.42) rectangle (6.95,0.65);
  % H2 is one connected plan-view sheet.  The white cut-outs are aligned
  % sensor windows; they are not separate hydrogel islands.
  \fill[AUBYellow!70] (8.0,0.42) rectangle (11.2,0.72);
  \draw[AUBYellow!90!black] (8.0,0.42) rectangle (11.2,0.72);
  \foreach \wx in {8.55,9.43,10.31}{
    \fill[white] (\wx,0.49) rectangle ++(0.38,0.16);
    \draw[AUBGray] (\wx,0.49) rectangle ++(0.38,0.16);
  }
  \fill[AUBCoral!55] (12.0,0.42) rectangle (15.2,0.65);
  \node[font=\fontsize{6.6}{7.5}\selectfont,text=AUBRed] at (1.6,0.57)
    {sensor plane};
\end{tikzpicture}
\end{center}
{\fontsize{7.5}{8.7}\selectfont\color{AUBGray}
\textbf{Figure Q2-H3. Evidence-informed hydrogel comparison.} H0 is the required
calibration baseline; H1 uses separate pads, H2 is one connected sheet with
aligned sensor windows (white cut-outs), and H3 covers the sensor regions. At
the humidity location, an H2 window would use a hydrophobic vapour-permeable
barrier rather than an open liquid path. The diagram is not to scale.\par}

\vspace{0.75em}
\hypertarget{q2-hydrogel-structures}{}
\qtwosubsection{2.6.2 Selected structures and decision gate}

\begingroup
\fontsize{8.75}{10.3}\selectfont
\renewcommand{\arraystretch}{1.14}
\begin{tabularx}{\linewidth}{>{\centering\arraybackslash\bfseries\color{AUBRed}}p{1.05cm}>{\raggedright\arraybackslash}p{3.3cm}Y Y}
\rowcolor{AUBRed}
\color{white}\textbf{Option} &
\color{white}\textbf{Structure} &
\color{white}\textbf{Purpose} &
\color{white}\textbf{Decision} \\
H0 & Bare printed PDMS with minimum protective geometry. &
Establish the reference calibration and the fastest possible transfer of each
stimulus to its sensor. & \textbf{Mandatory baseline.} Test first. \\
\rowcolor{AUBPale}
H1 & Patterned hydrogel islands outside active sensor footprints and conductor
exits. & Test skin feel and local moisture uptake while minimizing
direct interference with sensing zones. & \textbf{Recommended first hydrogel
experiment.} Compare directly with H0. \\
H2 & Connected windowed hydrogel sheet covering the spaces between sensors;
aligned cut-outs expose pressure, shear, and temperature footprints, while a
hydrophobic vapour-permeable vent protects the humidity route. & Explore broader
skin coverage without placing wet gel directly over active sensing zones. & Test
only if H1 passes; reject H2 if window edges redistribute load, trap liquid, or
shift calibration beyond limits. \\
\rowcolor{AUBPale}
H3 & Continuous hydrogel over all sensor regions. & Maximum continuous soft
contact, but maximum risk of buffering and cross-sensitivity. & \textbf{Not
recommended for the first build.} Activate only after H0--H2 evidence. \\
\end{tabularx}
\endgroup
\sourceanchor{Selected H0--H3 implementation sequence. A bonded
hydrogel/PDMS interface is literature-supported \cite{yuk2016}; its
effect on this multimodal slab remains experimental.}

\vspace{0.65em}
\evidencebox{Phase 3 pass condition:}{Select a hydrogel configuration only if it
has a repeatable PDMS bond, acceptable swelling and drying, no damaging liquid
path to the silver conductors, and predefined limits for humidity delay,
temperature lag, pressure attenuation, shear-vector error, hysteresis, drift,
and localization. Otherwise retain the bare calibrated PDMS interface.}

\clearpage
\sectiontitle{Q2. Whole-Liner Integration Pathway}

\hypertarget{q2-whole-liner}{}
\qtwosubsection{2.7 From the validated slab to a molded silicone liner}

{\fontsize{9.65}{11.4}\selectfont
The whole liner is the subsequent systems-integration stage. Printing and calibrating
the flat PDMS slab first isolates substrate and sensor problems before curvature,
pre-strain, cable routing, user fit, and repeated donning are introduced. The
complete liner will be assembled from validated sensor regions rather than
assuming that a successful flat print can be copied directly onto a curved
surface.\par}

\begingroup
\fontsize{8.8}{10.35}\selectfont
\renewcommand{\arraystretch}{1.14}
\begin{tabularx}{\linewidth}{>{\raggedright\arraybackslash\bfseries\color{AUBRed}}p{3.15cm}Y Y}
\rowcolor{AUBRed}
\color{white}\textbf{Integration route} &
\color{white}\textbf{Advantage} &
\color{white}\textbf{Main risk and gate} \\
Validated PDMS sensor tiles placed during molding & Preserves a tested printing
surface and allows defective tiles to be rejected before liner assembly. &
Tile edges, bonding, local stiffness, wiring exits, and strain transfer must be
tested under liner-like curvature and compression. \\
\rowcolor{AUBPale}
Printed flat PDMS sheet formed or laminated into the liner & Keeps printing on
a planar surface while enabling a larger connected array. & Forming can impose
pre-strain, wrinkle conductors, move coordinates, and change calibration. \\
Direct printing on a curved cured liner & Reduces assembly interfaces in
principle. & Most difficult first route: printhead clearance, drop landing,
wetting, focal distance, curing, and inspection vary with curvature. \\
\rowcolor{AUBPale}
Recommended sequence & Calibrate flat sensor slabs or tiles, test them under
known curvature, then integrate only passing units into a molded liner. &
Recalibrate after integration and compare with the original flat baseline. \\
\end{tabularx}
\endgroup
\sourceanchor{Integration routes selected for comparative implementation.
Published personalized-liner
work supports staged sensor integration and post-assembly testing, not any one
of these manufacturing routes \cite{paterno2020}.}

\vspace{0.35em}
\begin{center}
\begin{tikzpicture}[>=stealth,node distance=3mm]
  \tikzstyle{stage}=[draw=AUBLine,fill=AUBPale,rounded corners=1pt,
    minimum height=7mm,text width=33mm,align=center,
    font=\fontsize{7}{8}\selectfont]
  \node[stage] (s1) {flat PDMS\\calibration};
  \node[stage,right=of s1] (s2) {curvature and\\pre-strain check};
  \node[stage,right=of s2] (s3) {tile/liner\\integration};
  \node[stage,right=of s3] (s4) {post-molding\\recalibration};
  \draw[->,thick,color=AUBRed] (s1)--(s2);
  \draw[->,thick,color=AUBRed] (s2)--(s3);
  \draw[->,thick,color=AUBRed] (s3)--(s4);
\end{tikzpicture}
\end{center}
{\fontsize{7.5}{8.7}\selectfont\color{AUBGray}
\textbf{Figure Q2-L1. Whole-liner evidence ladder.} Coordinate accuracy and
calibration are rechecked whenever the mechanical boundary condition changes.\par}

\vspace{0.8em}
\begin{minipage}[t]{0.485\textwidth}
  \vspace{0pt}
  \figureplaceholder{54mm}{Dimensioned printed-slab section}{
  Add the final orthographic plan and cross-section. Label printed layers,
  active sensor footprints, protective PDMS, conductor exits, coordinate
  origin, nominal 5.2 mm thickness, and verified dimensions.}
\end{minipage}\hfill
\begin{minipage}[t]{0.485\textwidth}
  \vspace{0pt}
  \figureplaceholder{54mm}{Molded-liner integration design}{
  Show validated sensor slabs or tiles placed at defined liner coordinates,
  routing toward the acquisition module, curvature, molding boundaries, and
  the sequence of post-integration calibration tests.}
\end{minipage}

\vspace{0.65em}
\evidencebox{Stage boundary:}{The CAD liner render defines the integration-stage
design. It becomes works-like evidence after the flat printed PDMS slab
passes its electrical, sensor, environmental, mechanical, and coordinate gates.}

\clearpage
\sectiontitle{Q2. Literature Synthesis and Implementation Roadmap}

\hypertarget{q2-synthesis}{}
\qtwosubsection{2.8 Reused solutions, completed work, and integration focus}

\begingroup
\fontsize{8.5}{10}\selectfont
\renewcommand{\arraystretch}{1.1}
\begin{tabularx}{\linewidth}{>{\raggedright\arraybackslash\bfseries\color{AUBRed}}p{2.9cm}Y Y}
\rowcolor{AUBRed}
\color{white}\textbf{Literature-backed block} &
\color{white}\textbf{Project evidence available now} &
\color{white}\textbf{System-specific implementation output} \\
Printed sensor topologies & Pressure/shear, temperature, and humidity artwork
has been created and printed together on the large PET process sheet. & Transfer
the established geometries to PDMS and produce separate calibration,
cross-sensitivity, repeatability, and drift records for every channel. \\
\rowcolor{AUBPale}
Direct PDMS conductors & A 5.2 mm PDMS slab has been molded; two published
printing routes define starting controls for wetting and mechanical survival
\cite{albrecht2018stretch,sun2016}. & Record the project's surface treatment,
print window, profile, continuity, adhesion, curing, swelling, and cyclic
electrical response. \\
Coordinate-aware sensor layout & Fusion/AutoCAD slab renders and a placement
blueprint define relative sensor positions; the acquisition and mapping
algorithm is under development. & Export a dimensioned coordinate table, bind
each electrical channel to one sensor and position, and verify the mapping with
known-location stimuli. \\
\rowcolor{AUBPale}
Sensorized liner systems & Embedded and wireless prosthetic-interface studies
provide precedents for synchronized distributed sensing
\cite{paterno2020,q3kwak2020}. & Preserve channel identity through flat-slab,
curvature, molding, connection routing, and post-integration recalibration. \\
Skin-facing material system & Hydrogel bonding and sweat-collection studies
identify workable material methods and important thickness/time effects
\cite{yuk2016,lu2025}. & Compare bare PDMS with patterned structures only after
the bare slab is calibrated; retain the structure that preserves mechanical,
thermal, and humidity measurements. \\
Spatial interpretation and bench model & Published socket maps, movement
references, and mock limbs provide visualization and controlled-test precedents
\cite{q3turner2021,noll2018,quinlan2020,phillips2025}. & Quantify localization,
spatial blur, update stability, and modality disagreement on a repeatable phantom
before assigning a user-facing discomfort-risk interpretation. \\
\end{tabularx}
\endgroup
\sourceanchor{Synthesis of the functional-block literature in Sections
2.3--2.7 and the project artifacts documented in Section 2.2. The right-hand
column states implementation outputs rather than completed results.}

\vspace{0.65em}
\hypertarget{q2-roadmap}{}
\qtwosubsection{2.9 Research-backed implementation roadmap}

\begingroup
\fontsize{8.3}{9.75}\selectfont
\renewcommand{\arraystretch}{1.09}
\begin{tabularx}{\linewidth}{>{\centering\arraybackslash\bfseries\color{AUBRed}}p{0.8cm}>{\raggedright\arraybackslash\bfseries}p{3.25cm}Y>{\raggedright\arraybackslash}p{2.7cm}}
\rowcolor{AUBRed}
\color{white}\textbf{Stage} &
\color{white}\textbf{Build or study} &
\color{white}\textbf{Evidence transition} &
\color{white}\textbf{Status} \\
A & Functional-block and closest-solution map & Identify the established
solution used for each block and bound the integration gap. & \textbf{Completed
in this review} \\
\rowcolor{AUBPale}
B & Full artwork printed on PET & Preserve the print record, continuity results,
defects, routing, and connection access as process evidence. & \textbf{Completed
build} \\
C & Molded 5.2 mm PDMS slab and direct print & Establish the actual surface
treatment, printed geometry, swelling, adhesion, electrical integrity, and
yield on the intended substrate. & \textbf{Active build} \\
\rowcolor{AUBPale}
D & Separate sensor calibrations & Generate pressure/shear matrices,
temperature calibration, humidity response/recovery, strain effects, and
cross-sensitivity before signal fusion. & \textbf{Follows the stable PDMS print} \\
E & Multimodal slab on controlled phantom & Apply known single and combined
stimuli; measure repeatability, drift, localization, spatial blur, and incorrect
channel activation. & \textbf{Follows separate calibration} \\
\rowcolor{AUBPale}
F & Patterned skin-interface comparison & Compare H0 bare PDMS with H1 islands
and the H2 windowed sheet for moisture handling, transfer error, drying,
adhesion, and cycling. &
\textbf{Follows the bare baseline} \\
G & Curved liner and spatial software integration & Recheck coordinates and
calibration after curvature/molding, then connect the verified channels to the
threshold and spatial-map pipeline. & \textbf{System-integration stage} \\
\end{tabularx}
\endgroup
\sourceanchor{Roadmap ordered from completed artifacts to substrate transfer,
modality calibration, controlled multimodal testing, skin-interface comparison,
and liner/software integration. Phantom and localization methods are informed by
\cite{quinlan2020,phillips2025,q3turner2021,noll2018}.}

\vspace{0.55em}
\evidencebox{Q2 conclusion:}{The literature does not justify re-inventing or
claiming novelty for the individual sensor mechanisms. It supports a build
strategy that leverages established pressure/shear, temperature, humidity,
PDMS-printing, liner-sensing, and spatial-interpretation solutions. The work
specific to this project is the measured integration of those blocks on the
5.2 mm PDMS slab and then in a coordinate-preserving molded liner. The completed
PET print, molded PDMS specimen, sensor artwork, and CAD layouts are reported as
completed work; later roadmap stages are not presented as measured results.}

\clearpage
